\documentclass[12pt]{article}
\usepackage[utf8]{inputenc}
\usepackage[T1]{fontenc}
\usepackage{hyperref}
\title{Tangchangya Language Resources for K-12 Curriculum Development}
\author{Compiled by Claude Sonnet 5 (claude-sonnet-5) for Andr\'as Kornai}
\date{}
\begin{document}
\maketitle

\section{Introduction}
Tangchangya (ISO 639-3 tnv) is also spelled Tanchangya or Tongchangya. It is spoken by an indigenous community in the Chittagong Hill Tracts of Bangladesh. In Bangladesh it is written with Bengali letters, and it also has its own script, the Tanchangya script (Ka-Pat). One source says that script was introduced in Myanmar in 2012. The search found no primer or mother-tongue textbook for Tangchangya specifically.

\section{Primary Educational Resources}
Press coverage of the government's mother-tongue textbook initiative for pre-primary students names Chakma, Marma, Garo, Sadri and Tripura, but not Tanchangya. No Tanchangya primer or textbook turned up.
\begin{itemize}
\item Financial Express, indigenous textbooks in school: \url{https://old.thefinancialexpress.com.bd/public/print/indigenous-textbooks-in-school-1644681675}.
\item Dhaka Tribune editorial, ``Saving our languages'': \url{https://www.dhakatribune.com/opinion/editorial/12888/saving-our-languages}.
\end{itemize}

\section{Linguistic and Reference Materials}
\begin{itemize}
\item Wikipedia, ``Tanchangya language'': \url{https://en.wikipedia.org/wiki/Tanchangya_language}.
\item Wikipedia, ``Tanchangya script'': \url{https://en.wikipedia.org/wiki/Tanchangya_script}.
\item Omniglot, Tanchangya: \url{https://omniglot.com/writing/tanchangya.htm}.
\end{itemize}

\section{Religious and Community Resources}
Nothing specific to Tangchangya turned up in the search results.

\section{Digital and Audio Resources}
No audio, video or software resources turned up in the search results.

\section{Curriculum Development Context}
SIL Bangladesh notes the government's policy of supporting Mother Tongue Based Multilingual Education: \url{https://www.silbangladesh.org/node/1486}. Tanchangya is not named among the languages with textbooks so far. A curriculum effort would need to settle between Bengali-based and Tanchangya-script orthography and could look to the Chakma and Marma textbook programmes as a model. Next steps would be to check with the National Curriculum and Textbook Board and SIL Bangladesh.

\end{document}